\documentclass[onecolumn,aps,prb,amsmath,amssymb,superscriptaddress,longbibliography]{revtex4-2}

\usepackage{newtxtext}
\usepackage{newtxmath}
\usepackage{microtype}
\usepackage{bm}
\usepackage{mathtools}
\usepackage{booktabs}
\usepackage{array}
\usepackage{multirow}
\usepackage[table]{xcolor}
\usepackage{graphicx}
\usepackage{rotating}
\definecolor{PRLBlue}{RGB}{46,48,146}
\definecolor{SemigroupJ}{RGB}{218,238,222}
\definecolor{SemigroupI}{RGB}{218,234,244}
\definecolor{GroupOnly}{RGB}{239,241,242}
\definecolor{TRBrown}{RGB}{126,68,18}
\usepackage[
    colorlinks,
    citecolor=PRLBlue,
    linkcolor=PRLBlue,
    urlcolor=PRLBlue,
]{hyperref}
\usepackage{orcidlink}

\newcommand{\Tr}{\operatorname{Tr}}
\newcommand{\ee}{\mathrm{e}}
\newcommand{\ii}{\mathrm{i}}
\newcommand{\trans}{\mathsf{T}}
\newcommand{\mc}[1]{\mathcal{#1}}
\newcommand{\bb}[1]{\bm{#1}}
\newcommand{\ror}[2]{\href{https://ror.org/#1}{#2}}
\newcommand{\groupcell}{\cellcolor{GroupOnly}\shortstack{group only}}
\newcommand{\trgroupcell}{\cellcolor{GroupOnly}\shortstack{group only\\[-1pt]\textcolor{TRBrown}{\bfseries TR positive}}}
\newcommand{\jcell}{\cellcolor{SemigroupJ}\shortstack{semigroup\\[-1pt]\(\eta=J_2g\)}}
\newcommand{\icell}{\cellcolor{SemigroupI}\shortstack{semigroup\\[-1pt]\(\eta=\ii J_2g\)}}
\newcommand{\jonehead}[2]{\shortstack{\(J_1^\trans gJ_1=#1g\)\\\(J_1^2=#2I\)}}
\newcommand{\jtwohead}[3]{\shortstack{\(J_2^\trans gJ_2=#1g\)\\\(J_2^2=#2I\)\\\(J_1J_2=#3J_2J_1\)}}

\allowdisplaybreaks

\begin{document}

\title{Enumeration of Contraction Semigroups for Sign-Problem-Free Quantum Monte Carlo}

\author{Yin-Kai Yu\,\orcidlink{0009-0006-2589-2772}}
\affiliation{Beijing National Laboratory for Condensed Matter Physics\\ \& \ror{05cvf7v30}{Institute of Physics, Chinese Academy of Sciences}, Beijing 100190, China}
\affiliation{\ror{05qbk4x57}{University of Chinese Academy of Sciences}, Beijing 100049, China}

\author{Zhong-Chao Wei\,\orcidlink{0000-0002-1708-4306}}
\email{zcwei@thp.uni-koeln.de}
\affiliation{Qingdao Binhai University, Qingdao 266555, China}

\author{Lei Wang\,\orcidlink{0009-0004-1348-8583}}
\email{wanglei@iphy.ac.cn}
\affiliation{Beijing National Laboratory for Condensed Matter Physics\\ \& \ror{05cvf7v30}{Institute of Physics, Chinese Academy of Sciences}, Beijing 100190, China}
\affiliation{\ror{05qbk4x57}{University of Chinese Academy of Sciences}, Beijing 100049, China}
\affiliation{\ror{020vtf184}{Songshan Lake Materials Laboratory}, Dongguan, Guangdong 523808, China}

\author{Zi-Xiang Li\,\orcidlink{0000-0002-3941-365X}}
\email{zixiangli@iphy.ac.cn}
\affiliation{Beijing National Laboratory for Condensed Matter Physics\\ \& \ror{05cvf7v30}{Institute of Physics, Chinese Academy of Sciences}, Beijing 100190, China}
\affiliation{\ror{05qbk4x57}{University of Chinese Academy of Sciences}, Beijing 100049, China}

\date{\today}

\begin{abstract}
Auxiliary-field quantum Monte Carlo expresses each field configuration as an ordered product of Gaussian many-body operators.  A contraction semigroup keeps this product inside a connected region where a real weight cannot cross a determinant zero or a bosonic pole.  We formulate the construction in a common Nambu representation for bosons and fermions and classify two real orthogonal matrices through five binary relations.  Of the \(32\) bosonic and \(32\) fermionic classes, exactly eight in each set admit a contraction cone with negative-definite generators.  Four use the metric \(J_2g\), four use \(\ii J_2g\), and the remaining \(24\) contain only the maximal subgroup.  Two short algebraic criteria give the complete classification.  We enumerate the maximal-subgroup and contraction-cone generators for every allowed class and translate them into number-conserving and anomalous Hamiltonian terms.  Time-reversal positivity supplies an independent sign-free route for eight bosonic and four fermionic group-only classes.  For trace-convergent realizations connected to a positive reference, the resulting table gives a constructive map from symmetry algebra to sign-problem-free auxiliary-field decompositions.
\end{abstract}

\maketitle

\tableofcontents

\section{Quadratic evolution and Monte Carlo weights}

A Hubbard--Stratonovich decomposition turns an interacting partition function into an integral or sum over auxiliary fields~\cite{blankenbecler1981prd}.  Once the fields are fixed, the quantum particles evolve through quadratic operators
\begin{equation}
    p=\Tr\!\left(\ee^{\hat h_1}\ee^{\hat h_2}\cdots\ee^{\hat h_L}\right),
    \qquad
    \hat h_\ell=\frac12\hat\chi^\trans gA_\ell\hat\chi .
    \label{eq:configuration_weight}
\end{equation}
The Nambu column
\begin{equation}
    \hat\chi=(\hat a_1,\ldots,\hat a_N,
    \hat a_1^+,\ldots,\hat a_N^+)^\trans
    \label{eq:nambu_column}
\end{equation}
contains annihilation and creation operators.  We use \(\hat a\) for either bosons or fermions in the common formulas and distinguish them when an operator identity depends on statistics.  Their canonical matrices are
\begin{equation}
    g_{\rm B}=\begin{pmatrix}0&-I_N\\ I_N&0\end{pmatrix},
    \qquad
    g_{\rm F}=\begin{pmatrix}0&I_N\\ I_N&0\end{pmatrix}.
    \label{eq:canonical_g}
\end{equation}
The condition
\begin{equation}
    gA_\ell+A_\ell^\trans g=0
    \label{eq:quadratic_lie_algebra}
\end{equation}
makes \(gA_\ell\) symmetric for bosons and antisymmetric for fermions, as required by their canonical algebra.

The ordered one-body propagation matrix is
\begin{equation}
    M=\ee^{A_1}\ee^{A_2}\cdots\ee^{A_L}.
    \label{eq:propagation_matrix}
\end{equation}
The corresponding Gaussian trace formulas are~\cite{munn1978ensemble,klich2014jstat}
\begin{equation}
    p_{\rm B}=\left[(-1)^N\det(I_{2N}-M)\right]^{-1/2},
    \qquad
    p_{\rm F}=\left[\det(I_{2N}+M)\right]^{1/2}.
    \label{eq:gaussian_weights}
\end{equation}
The square roots retain the sign selected by the ordered Fock-space product in Eq.~\eqref{eq:configuration_weight}.  A sign-free construction therefore needs a real weight, a connected path from a positive reference configuration, and no zero or pole along that path.

The contraction-semigroup approach introduced by Wei realizes these requirements directly~\cite{wei2024prb}.  Let \(\eta=\eta^\dagger\) and \(\eta^2=I\).  A slice generator obeying
\begin{equation}
    \eta A_\ell+A_\ell^\dagger\eta\leq0
    \label{eq:generator_contraction}
\end{equation}
produces
\begin{equation}
    \ee^{A_\ell^\dagger}\eta\ee^{A_\ell}\leq\eta .
    \label{eq:matrix_contraction}
\end{equation}
Products preserve the same inequality because
\begin{align}
    &(M_1M_2)^\dagger\eta(M_1M_2)-\eta \notag\\
    &\quad=M_2^\dagger(M_1^\dagger\eta M_1-\eta)M_2
      +(M_2^\dagger\eta M_2-\eta).
    \label{eq:semigroup_product}
\end{align}
If one slice obeys Eq.~\eqref{eq:matrix_contraction} with \(<\) in place of \(\leq\), and all other slices are invertible contractions, Eq.~\eqref{eq:semigroup_product} gives the same strict inequality for the complete product.  Such a product has no eigenvalue on the unit circle.  In particular, \(+1\) cannot generate a bosonic pole and \(-1\) cannot generate a fermionic determinant zero.  When an anti-linear constraint makes the determinant in Eq.~\eqref{eq:gaussian_weights} real, it remains on the sign selected at a positive reference configuration.  The ordered square root is continuous along the same path, so the Monte Carlo weight cannot change sign.  Boundary configurations follow as limits of strict contractions.

For a bosonic grand-canonical simulation, the quadratic evolution must also damp high-occupation states so that the many-body trace converges.  Strict propagation in the standard particle--hole realization with \(\eta=\operatorname{diag}(I_N,-I_N)\) has this property.  With a more general metric, one checks directly that the quadratic evolution makes the occupation-number sum converge.

\section{Two matrices and five relations}

Both \(J_1\) and \(J_2\) act on the Nambu column and are real orthogonal matrices, \(J_j^\trans J_j=I\).
The first matrix, \(J_1\), supplies the anti-linear constraint
\begin{equation}
    J_1^\trans A_\ell J_1=\overline{A_\ell}.
    \label{eq:J1_reality}
\end{equation}
It follows slice by slice that \(J_1^\trans M J_1=\overline M\).  Hence the determinants in Eq.~\eqref{eq:gaussian_weights} are real.  Together with strict contraction and a positive reference configuration, this is the reality condition used by the semigroup proof above.

The second real orthogonal matrix, \(J_2\), builds the Hermitian metric used in Eq.~\eqref{eq:generator_contraction}.  The discrete information needed from \(J_1,J_2\) is
\begin{equation}
\begin{aligned}
    J_1^\trans gJ_1&=s_1g,&\qquad J_1^2&=q_1I,\\
    J_2^\trans gJ_2&=s_2g,& J_2^2&=q_2I,\\
    J_1J_2&=rJ_2J_1,
\end{aligned}
\qquad
    s_1,q_1,s_2,q_2,r\in\{+1,-1\}.
    \label{eq:five_relations}
\end{equation}
The first four signs say whether each matrix preserves or reverses \(g\), and whether it squares to \(+I\) or \(-I\).  The fifth sign records whether the two transformations commute or anticommute.  These are the five independent binary choices, giving \(2^5=32\) classes for bosons and another \(32\) for fermions.

The signs carried by \(J_2\) determine whether \(J_2g\) is Hermitian or anti-Hermitian.  The appropriate metric is listed in Table~\ref{tab:metric_choice}.  In every row it satisfies \(\eta^\dagger=\eta\) and \(\eta^2=I\).
\begin{table}[tb]
\caption{Metric selected by the two \(J_2\) relations.  Multiplication by \(\ii\) converts an anti-Hermitian \(J_2g\), whose square is \(-I\), into a Hermitian matrix satisfying \(\eta^2=I\).}
\label{tab:metric_choice}
\centering
\begin{ruledtabular}
\begin{tabular}{cccc}
Statistics & \(J_2^\trans gJ_2\) & \(J_2^2\) & \(\eta\)\\
\hline
boson & \(-g\) & \(+I\) & \(J_2g\)\\
boson & \(+g\) & \(+I\) & \(\ii J_2g\)\\
boson & \(-g\) & \(-I\) & \(\ii J_2g\)\\
boson & \(+g\) & \(-I\) & \(J_2g\)\\
fermion & \(-g\) & \(-I\) & \(J_2g\)\\
fermion & \(+g\) & \(-I\) & \(\ii J_2g\)\\
fermion & \(-g\) & \(+I\) & \(\ii J_2g\)\\
fermion & \(+g\) & \(+I\) & \(J_2g\)
\end{tabular}
\end{ruledtabular}
\end{table}

Complex conjugation of Eq.~\eqref{eq:quadratic_lie_algebra} gives \(A^\dagger g=-g\overline A\).  Substitution of the metric from Table~\ref{tab:metric_choice} then gives the identity
\begin{equation}
    \eta A^\dagger\eta=-J_2^\trans\overline A J_2
    \label{eq:eta_J2_identity}
\end{equation}
which separates \(A\) into two physically distinct parts,
\begin{align}
    A_{\rm group}
      &=\frac12\left(A+J_2^\trans\overline A J_2\right),
    &J_2^\trans A_{\rm group}J_2&=\overline{A_{\rm group}},
    \label{eq:A_group_projection}\\
    A_{\rm cone}
      &=\frac12\left(A-J_2^\trans\overline A J_2\right),
    &J_2^\trans A_{\rm cone}J_2&=-\overline{A_{\rm cone}}.
    \label{eq:A_cone_projection}
\end{align}
The first part generates the maximal subgroup that preserves \(\eta\).  It contributes zero to the contraction matrix.  The second part points into or out of the contraction region.  Its allowed direction is fixed by
\begin{equation}
    \eta A+A^\dagger\eta=2\eta A_{\rm cone}\leq0.
    \label{eq:cone_condition}
\end{equation}
Thus \(A_{\rm group}\) describes reversible motion that preserves the metric, while a negative-definite \(\eta A_{\rm cone}\) supplies damping.  Auxiliary fields that take both signs naturally couple to group directions.  They may also enter cone directions when the complete field-dependent combination remains in the cone.

\section{Enumeration and periodic table}

The enumeration uses the real Pauli basis
\begin{equation}
    I=I_2,\qquad X=\sigma_x,\qquad
    Y=\ii\sigma_y=\begin{pmatrix}0&1\\-1&0\end{pmatrix},
    \qquad Z=\sigma_z .
    \label{eq:real_pauli_basis}
\end{equation}
Tensor products of these four matrices generate real orthogonal representatives.  For example, \(IYX\) means \(I_2\otimes Y\otimes X\), and the first factor is always the Nambu factor.  A GNU Octave calculation searched products with one, two, and three Pauli factors, retained the smallest representative of every sign class, and solved the linear equations
\begin{equation}
\begin{gathered}
    gA+A^\trans g=0,\qquad
    J_1^\trans AJ_1=\overline A,\\
    J_2^\trans A_{\rm group}J_2=\overline{A_{\rm group}},\qquad
    J_2^\trans A_{\rm cone}J_2=-\overline{A_{\rm cone}}.
\end{gathered}
    \label{eq:enumeration_equations}
\end{equation}
Complex matrices were separated into real and imaginary parts, so Eq.~\eqref{eq:enumeration_equations} became an ordinary real linear system for the matrix entries.  Solving that system produced every allowed coefficient.  The final step tested whether the cone subspace intersects \(\eta A_{\rm cone}<0\).  All \(64\) classes were evaluated, and every displayed Pauli direction was substituted back into the original matrix equations.

Table~\ref{tab:periodic_table} is the resulting periodic table.  The bosonic and fermionic sectors each have eight contraction-semigroup classes, equally divided between the two metrics, and \(24\) group-only classes.  A time-reversal marker identifies a group-only cell that admits a physical Majorana or Kramers representative.  The representatives are given in Sec.~\ref{sec:group_only_TR}.

\begin{sidewaystable}
\caption{Complete enumeration of the five relations in Eq.~\eqref{eq:five_relations}.  Rows specify the action and square of \(J_1\).  Columns specify the action and square of \(J_2\), followed by their mutual sign.  Green and blue cells contain a contraction cone with a negative-definite direction.  Gray cells contain only the maximal subgroup.  The label ``TR positive'' marks a group-only class with at least one physical representative whose two induced time reversals form a Majorana \(\{+,-\}\) or Kramers \(\{-,-\}\) pair.}
\label{tab:periodic_table}
\centering
\setlength{\tabcolsep}{3.5pt}
\renewcommand{\arraystretch}{1.55}
\scriptsize
\resizebox{0.98\textheight}{!}{%
\begin{tabular}{c*{8}{c}}
\toprule
\multicolumn{9}{c}{\bfseries Bosons}\\
\cmidrule(lr){1-9}
& \jtwohead{+}{+}{+}
& \jtwohead{+}{+}{-}
& \jtwohead{+}{-}{+}
& \jtwohead{+}{-}{-}
& \jtwohead{-}{+}{+}
& \jtwohead{-}{+}{-}
& \jtwohead{-}{-}{+}
& \jtwohead{-}{-}{-}\\
\midrule
\jonehead{+}{+}
&\groupcell&\icell&\groupcell&\trgroupcell&\jcell&\groupcell&\trgroupcell&\groupcell\\
\jonehead{+}{-}
&\groupcell&\icell&\groupcell&\trgroupcell&\jcell&\groupcell&\trgroupcell&\groupcell\\
\jonehead{-}{+}
&\icell&\groupcell&\trgroupcell&\groupcell&\groupcell&\jcell&\groupcell&\trgroupcell\\
\jonehead{-}{-}
&\icell&\groupcell&\trgroupcell&\groupcell&\groupcell&\jcell&\groupcell&\trgroupcell\\
\midrule
\multicolumn{9}{c}{\bfseries Fermions}\\
\cmidrule(lr){1-9}
& \jtwohead{+}{+}{+}
& \jtwohead{+}{+}{-}
& \jtwohead{+}{-}{+}
& \jtwohead{+}{-}{-}
& \jtwohead{-}{+}{+}
& \jtwohead{-}{+}{-}
& \jtwohead{-}{-}{+}
& \jtwohead{-}{-}{-}\\
\midrule
\jonehead{+}{+}
&\groupcell&\groupcell&\groupcell&\icell&\groupcell&\groupcell&\jcell&\groupcell\\
\jonehead{+}{-}
&\groupcell&\trgroupcell&\groupcell&\icell&\trgroupcell&\groupcell&\jcell&\groupcell\\
\jonehead{-}{+}
&\groupcell&\groupcell&\icell&\groupcell&\groupcell&\groupcell&\groupcell&\jcell\\
\jonehead{-}{-}
&\trgroupcell&\groupcell&\icell&\groupcell&\groupcell&\trgroupcell&\groupcell&\jcell\\
\bottomrule
\end{tabular}}
\end{sidewaystable}

\section{Two screening lemmas}

The \(64\)-cell pattern follows from two short calculations.  They also show that the enumeration is complete within the five relations.

\paragraph{Lemma 1: the square of \(J_2\).}
For bosons, \(g^\trans=-g\) and Eq.~\eqref{eq:quadratic_lie_algebra} gives \((gA)^\trans=gA\).  If \(J_2^2=-I\), orthogonality implies \(J_2^\trans=-J_2\), and hence
\begin{equation}
\begin{aligned}
    \Tr(J_2gA)
    &=\Tr[(J_2gA)^\trans]
      =\Tr[(gA)^\trans J_2^\trans]\\
    &=-\Tr(gAJ_2)=-\Tr(J_2gA)=0.
\end{aligned}
    \label{eq:boson_trace_obstruction}
\end{equation}
The same result holds after multiplication by \(\ii\).  A negative-semidefinite Hermitian matrix with zero trace must vanish.  A bosonic contraction cone with a negative-definite direction therefore requires
\begin{equation}
    J_2^2=+I.
    \label{eq:boson_J2_square}
\end{equation}

For fermions, \(g^\trans=g\) gives \((gA)^\trans=-gA\).  If \(J_2^2=+I\), then \(J_2^\trans=J_2\), and
\begin{equation}
    \Tr(J_2gA)
    =\Tr[(gA)^\trans J_2^\trans]
    =-\Tr(J_2gA)=0.
    \label{eq:fermion_trace_obstruction}
\end{equation}
The fermionic requirement is consequently
\begin{equation}
    J_2^2=-I.
    \label{eq:fermion_J2_square}
\end{equation}
The opposite squares in Eqs.~\eqref{eq:boson_J2_square} and \eqref{eq:fermion_J2_square} arise directly from the symmetric bosonic \(gA\) and antisymmetric fermionic \(gA\).

\paragraph{Lemma 2: compatibility with the real-weight condition.}
Let \(H=\eta A_{\rm cone}=H^\dagger\), and write \(\eta=cJ_2g\), where Table~\ref{tab:metric_choice} fixes \(c=1\) or \(c=\ii\).  Once the square condition in Lemma~1 holds, the same table gives \(c/\overline c=-s_2\).  The transformation of \(H\) is then
\begin{equation}
\begin{aligned}
J_1^\trans HJ_1
&=c(J_1^\trans J_2J_1)(J_1^\trans gJ_1)
  (J_1^\trans A_{\rm cone}J_1)\\
&=crs_1J_2g\overline{A_{\rm cone}}
 =-rs_1s_2\overline H.
\end{aligned}
    \label{eq:H_transformation}
\end{equation}
The right-hand side equals \(\overline H\) precisely when
\begin{equation}
    r=-s_1s_2.
    \label{eq:commutation_rule}
\end{equation}
This equation has a direct reading.  If \(J_1\) and \(J_2\) both preserve \(g\), or both reverse it, they must anticommute.  If one preserves \(g\) and the other reverses it, they must commute.

For the other mutual relation, Eq.~\eqref{eq:H_transformation} reads \(J_1^\trans HJ_1=-\overline H\).  Starting from \(H\leq0\), one obtains
\begin{equation}
    J_1^\trans HJ_1\leq0
    \quad\Longrightarrow\quad
    -\overline H\leq0
    \quad\Longrightarrow\quad
    H\geq0.
    \label{eq:cone_sign_flip}
\end{equation}
The two inequalities force
\begin{equation}
    H=0.
    \label{eq:reality_flip_boundary}
\end{equation}
This leaves only the maximal-subgroup boundary.

The two requirements are also sufficient.  Choose \(\eta\) from Table~\ref{tab:metric_choice} and set
\begin{equation}
    A_{\rm cone}=-\eta .
    \label{eq:universal_cone_witness}
\end{equation}
Equations~\eqref{eq:boson_J2_square} or \eqref{eq:fermion_J2_square} ensure \(gA_{\rm cone}+A_{\rm cone}^\trans g=0\).  Equation~\eqref{eq:commutation_rule} ensures \(J_1^\trans A_{\rm cone}J_1=\overline{A_{\rm cone}}\), and the \(J_2\)-odd relation follows from the metric construction.  Finally,
\begin{equation}
    \eta A_{\rm cone}+A_{\rm cone}^\dagger\eta=-2I<0.
    \label{eq:universal_strict_contraction}
\end{equation}
We have therefore proved the complete criterion
\begin{equation}
\boxed{
\begin{array}{ll}
    J_2^2=+I,\quad r=-s_1s_2,&\text{bosons},\\[2pt]
    J_2^2=-I,\quad r=-s_1s_2,&\text{fermions}.
\end{array}}
    \label{eq:complete_semigroup_criterion}
\end{equation}
The sign \(q_1=\pm1\) remains free.  For bosons and fermions separately, \(s_1,q_1,s_2\) supply \(2^3=8\) allowed classes.  The four with \(s_2=-1\) use \(J_2g\), while the four with \(s_2=+1\) use \(\ii J_2g\).

\section{Generators and Hamiltonian terms}

The following tables give every generator for the smallest Pauli representative of each allowed class.  If the displayed representative has dimension \(d\), it can be tensored with an \(m\)-dimensional orbital, site, or flavor space.  Every solution has the form
\begin{equation}
    A_{\rm group}=\sum_P P\otimes K_P,
    \qquad
    A_{\rm cone}=\sum_P P\otimes K_P.
    \label{eq:pauli_generator_expansion}
\end{equation}
The four coefficient columns state the complete transpose and complex-conjugation conditions on \(K_P\).  The contraction condition acts on the full sum,
\begin{equation}
    \eta A_{\rm cone}
    =\eta\sum_P P\otimes K_P<0.
    \label{eq:full_coefficient_inequality}
\end{equation}
Explicitly, the four column headings mean
\begin{equation}
\begin{array}{c|cc}
&\overline{K_P}&K_P^\trans\\
\hline
\text{real symmetric}&K_P&K_P\\
\text{real antisymmetric}&K_P&-K_P\\
\text{imaginary symmetric}&-K_P&K_P\\
\text{imaginary antisymmetric}&-K_P&-K_P
\end{array}.
\label{eq:coefficient_types}
\end{equation}
It generally couples different \(K_P\) and cannot be replaced by separate sign conditions on the coefficients.  Each representative also lists the simple ray obtained by taking a real symmetric \(H>0\), setting the stated \(K_P\), and setting all other cone coefficients to zero.  Along that ray,
\begin{equation}
    A_{\rm cone}=-\eta\otimes H,
    \qquad
    \eta A_{\rm cone}=-I_d\otimes H<0.
    \label{eq:negative_ray}
\end{equation}

The first Pauli factor immediately identifies the physical quadratic term.  Let \(P=P_{\rm N}\otimes P_{\rm int}\) and \(L=P_{\rm int}\otimes K_P\).  Table~\ref{tab:operator_dictionary} translates each Nambu factor into operators.  The required symmetry of \(L\) follows automatically from the coefficient columns in Tables~\ref{tab:boson_ij2g_generators}--\ref{tab:fermion_j2g_generators}.
\begin{table}[tb]
\caption{Hamiltonian term generated by the first Pauli factor in \(A=P_{\rm N}\otimes L\).  The remaining Pauli factors and \(K_P\) are combined into \(L\).  Indices include every internal, orbital, and spatial label.}
\label{tab:operator_dictionary}
\centering
\small
\begin{ruledtabular}
\begin{tabular}{ccll}
\(P_{\rm N}\) & symmetry of \(L\) & bosonic term & fermionic term\\
\hline
\(I\) & \(L^\trans=-L\) &
\(\sum_{ij}b_i^+L_{ij}b_j\) &
\(\sum_{ij}c_i^+L_{ij}c_j\)\\[2pt]
\(Z\) & \(L^\trans=L\) &
\(\sum_{ij}b_i^+L_{ij}b_j+\frac12\Tr L\) &
\(\sum_{ij}c_i^+L_{ij}c_j-\frac12\Tr L\)\\[2pt]
\(X\) & \shortstack{\(L^\trans=L\), boson\\\(L^\trans=-L\), fermion} &
\(\frac12\sum_{ij}L_{ij}(b_i^+b_j^+-b_ib_j)\) &
\(\frac12\sum_{ij}L_{ij}(c_ic_j+c_i^+c_j^+)\)\\[2pt]
\(Y\) & \shortstack{\(L^\trans=L\), boson\\\(L^\trans=-L\), fermion} &
\(\frac12\sum_{ij}L_{ij}(b_i^+b_j^++b_ib_j)\) &
\(\frac12\sum_{ij}L_{ij}(-c_ic_j+c_i^+c_j^+)\)
\end{tabular}
\end{ruledtabular}
\end{table}

\begin{sidewaystable}
\caption{Bosonic contraction-semigroup generators with \(J_2^\trans gJ_2=g\), \(J_2^2=I\), and \(\eta=\ii J_2g\).  A Pauli word in a coefficient column means that it appears in Eq.~\eqref{eq:pauli_generator_expansion} with a matrix \(K_P\) of the stated type.  The tuple is \((s_1,q_1;s_2,q_2;r)\).  The complete negative-definite condition is Eq.~\eqref{eq:full_coefficient_inequality}.}
\label{tab:boson_ij2g_generators}
\centering
\renewcommand{\arraystretch}{2.2}
\footnotesize
\resizebox{0.98\textheight}{!}{%
\begin{tabular}{lcllll}
\toprule
Class and representative & generator & \shortstack{real\\symmetric} & \shortstack{real\\antisymmetric} & \shortstack{imaginary\\symmetric} & \shortstack{imaginary\\antisymmetric}\\
\midrule
\multirow{2}{0.22\textheight}{
\((+,+;+,+;-)\), \(d=4\)\\
\(J_1=IX, J_2=IZ\)\\
\(\eta=-\ii YZ\)\\
ray: \(K_{YZ}=\ii H\)}
& \(A_{\rm group}\) & XI, YI, ZI & II & IY & XY, YY, ZY\\
& \(A_{\rm cone}\) & XX, YX, ZX & IX & XZ, YZ, ZZ & IZ\\
\midrule
\multirow{2}{0.22\textheight}{
\((+,-;+,+;-)\), \(d=4\)\\
\(J_1=IY, J_2=IX\)\\
\(\eta=-\ii YX\)\\
ray: \(K_{YX}=\ii H\)}
& \(A_{\rm group}\) & XI, YI, ZI & II & XZ, YZ, ZZ & IZ\\
& \(A_{\rm cone}\) & IY & XY, YY, ZY & XX, YX, ZX & IX\\
\midrule
\multirow{2}{0.22\textheight}{
\((-,+;+,+;+)\), \(d=2\)\\
\(J_1=X, J_2=I\)\\
\(\eta=-\ii Y\)\\
ray: \(K_Y=\ii H\)}
& \(A_{\rm group}\) & X & I & -- & --\\
& \(A_{\rm cone}\) & -- & -- & Y, Z & --\\
\midrule
\multirow{2}{0.22\textheight}{
\((-,-;+,+;+)\), \(d=4\)\\
\(J_1=XY, J_2=II\)\\
\(\eta=-\ii YI\)\\
ray: \(K_{YI}=\ii H\)}
& \(A_{\rm group}\) & IY, XI, YX, YZ, ZX, ZZ & II, XY & -- & --\\
& \(A_{\rm cone}\) & -- & -- & XX, XZ, YI, ZI & IX, IZ, YY, ZY\\
\bottomrule
\end{tabular}}
\end{sidewaystable}

\begin{sidewaystable}
\caption{Bosonic contraction-semigroup generators with \(J_2^\trans gJ_2=-g\), \(J_2^2=I\), and \(\eta=J_2g\).  Notation follows Table~\ref{tab:boson_ij2g_generators}.}
\label{tab:boson_j2g_generators}
\centering
\renewcommand{\arraystretch}{2.2}
\footnotesize
\resizebox{0.98\textheight}{!}{%
\begin{tabular}{lcllll}
\toprule
Class and representative & generator & \shortstack{real\\symmetric} & \shortstack{real\\antisymmetric} & \shortstack{imaginary\\symmetric} & \shortstack{imaginary\\antisymmetric}\\
\midrule
\multirow{2}{0.22\textheight}{
\((+,+;-,+;+)\), \(d=2\)\\
\(J_1=I, J_2=X\)\\
\(\eta=Z\)\\
ray: \(K_Z=-H\)}
& \(A_{\rm group}\) & X & I & -- & --\\
& \(A_{\rm cone}\) & Y, Z & -- & -- & --\\
\midrule
\multirow{2}{0.22\textheight}{
\((+,-;-,+;+)\), \(d=4\)\\
\(J_1=IY, J_2=XI\)\\
\(\eta=ZI\)\\
ray: \(K_{ZI}=-H\)}
& \(A_{\rm group}\) & IY, XI & II, XY & YX, YZ, ZX, ZZ & --\\
& \(A_{\rm cone}\) & YI, ZI & YY, ZY & XX, XZ & IX, IZ\\
\midrule
\multirow{2}{0.22\textheight}{
\((-,+;-,+;-)\), \(d=2\)\\
\(J_1=X, J_2=Z\)\\
\(\eta=-X\)\\
ray: \(K_X=H\)}
& \(A_{\rm group}\) & -- & I & Y & --\\
& \(A_{\rm cone}\) & X & -- & Z & --\\
\midrule
\multirow{2}{0.22\textheight}{
\((-,-;-,+;-)\), \(d=4\)\\
\(J_1=XY, J_2=XX\)\\
\(\eta=ZX\)\\
ray: \(K_{ZX}=-H\)}
& \(A_{\rm group}\) & XI, YZ, ZZ & II & XZ, YI, ZI & IZ\\
& \(A_{\rm cone}\) & IY, YX, ZX & XY & XX & IX, YY, ZY\\
\bottomrule
\end{tabular}}
\end{sidewaystable}

\begin{sidewaystable}
\caption{Fermionic contraction-semigroup generators with \(J_2^\trans gJ_2=g\), \(J_2^2=-I\), and \(\eta=\ii J_2g\).  Notation follows Table~\ref{tab:boson_ij2g_generators}.}
\label{tab:fermion_ij2g_generators}
\centering
\renewcommand{\arraystretch}{2.2}
\scriptsize
\resizebox{0.98\textheight}{!}{%
\begin{tabular}{lcllll}
\toprule
Class and representative & generator & \shortstack{real\\symmetric} & \shortstack{real\\antisymmetric} & \shortstack{imaginary\\symmetric} & \shortstack{imaginary\\antisymmetric}\\
\midrule
\multirow{2}{0.22\textheight}{
\((+,+;+,-;-)\), \(d=4\)\\
\(J_1=IX, J_2=IY\)\\
\(\eta=\ii XY\)\\
ray: \(K_{XY}=-\ii H\)}
& \(A_{\rm group}\) & ZI & II, XI, YI & ZZ & IZ, XZ, YZ\\
& \(A_{\rm cone}\) & ZX & IX, XX, YX & IY, XY, YY & ZY\\
\midrule
\multirow{2}{0.22\textheight}{
\((+,-;+,-;-)\), \(d=8\)\\
\(J_1=IIY, J_2=IYX\)\\
\(\eta=\ii XYX\)\\
ray: \(K_{XYX}=-\ii H\)}
& \(A_{\rm group}\) &
\shortstack[l]{IXY, IYI, IZY, XXY, XYI\\XZY, YXY, YYI, YZY, ZII} &
\shortstack[l]{III, XII, YII\\ZXY, ZYI, ZZY} &
\shortstack[l]{IYZ, XYZ, YYZ\\ZIZ, ZXX, ZZX} &
\shortstack[l]{IIZ, IXX, IZX, XIZ, XXX\\XZX, YIZ, YXX, YZX, ZYZ}\\
& \(A_{\rm cone}\) &
\shortstack[l]{IIY, XIY, YIY\\ZXI, ZYY, ZZI} &
\shortstack[l]{IXI, IYY, IZI, XXI, XYY\\XZI, YXI, YYY, YZI, ZIY} &
\shortstack[l]{IYX, XYX, YYX\\ZIX, ZXZ, ZZZ} &
\shortstack[l]{IIX, IXZ, IZZ, XIX, XXZ\\XZZ, YIX, YXZ, YZZ, ZYX}\\
\midrule
\multirow{2}{0.22\textheight}{
\((-,+;+,-;+)\), \(d=4\)\\
\(J_1=YY, J_2=IY\)\\
\(\eta=\ii XY\)\\
ray: \(K_{XY}=-\ii H\)}
& \(A_{\rm group}\) & IY, YY & II, YI & -- & IX, IZ, YX, YZ\\
& \(A_{\rm cone}\) & ZX, ZZ & XX, XZ & XY, ZI & XI, ZY\\
\midrule
\multirow{2}{0.22\textheight}{
\((-,-;+,-;+)\), \(d=4\)\\
\(J_1=YI, J_2=IY\)\\
\(\eta=\ii XY\)\\
ray: \(K_{XY}=-\ii H\)}
& \(A_{\rm group}\) & IY, YY & II, YI & ZX, ZZ & XX, XZ\\
& \(A_{\rm cone}\) & -- & IX, IZ, YX, YZ & XY, ZI & XI, ZY\\
\bottomrule
\end{tabular}}
\end{sidewaystable}

\begin{sidewaystable}
\caption{Fermionic contraction-semigroup generators with \(J_2^\trans gJ_2=-g\), \(J_2^2=-I\), and \(\eta=J_2g\).  Notation follows Table~\ref{tab:boson_ij2g_generators}.}
\label{tab:fermion_j2g_generators}
\centering
\renewcommand{\arraystretch}{2.2}
\footnotesize
\resizebox{0.98\textheight}{!}{%
\begin{tabular}{lcllll}
\toprule
Class and representative & generator & \shortstack{real\\symmetric} & \shortstack{real\\antisymmetric} & \shortstack{imaginary\\symmetric} & \shortstack{imaginary\\antisymmetric}\\
\midrule
\multirow{2}{0.22\textheight}{
\((+,+;-,-;+)\), \(d=2\)\\
\(J_1=I, J_2=Y\)\\
\(\eta=Z\)\\
ray: \(K_Z=-H\)}
& \(A_{\rm group}\) & -- & I, Y & -- & --\\
& \(A_{\rm cone}\) & Z & X & -- & --\\
\midrule
\multirow{2}{0.22\textheight}{
\((+,-;-,-;+)\), \(d=4\)\\
\(J_1=IY, J_2=YI\)\\
\(\eta=ZI\)\\
ray: \(K_{ZI}=-H\)}
& \(A_{\rm group}\) & IY, YY & II, YI & ZX, ZZ & XX, XZ\\
& \(A_{\rm cone}\) & XY, ZI & XI, ZY & -- & IX, IZ, YX, YZ\\
\midrule
\multirow{2}{0.22\textheight}{
\((-,+;-,-;-)\), \(d=2\)\\
\(J_1=Z, J_2=Y\)\\
\(\eta=Z\)\\
ray: \(K_Z=-H\)}
& \(A_{\rm group}\) & -- & I & -- & X\\
& \(A_{\rm cone}\) & Z & -- & -- & Y\\
\midrule
\multirow{2}{0.22\textheight}{
\((-,-;-,-;-)\), \(d=4\)\\
\(J_1=YI, J_2=ZY\)\\
\(\eta=YY\)\\
ray: \(K_{YY}=-H\)}
& \(A_{\rm group}\) & IY & II, YX, YZ & XY, ZX, ZZ & XI\\
& \(A_{\rm cone}\) & YY & IX, IZ, YI & ZI & XX, XZ, ZY\\
\bottomrule
\end{tabular}}
\end{sidewaystable}

All representatives in the tables have Nambu dimension two or four, except the fermionic class \((+,-;+,-;-)\).  In that class both \(J_1\) and \(J_2\) preserve \(g\), square to \(-I\), and anticommute.  The \(+1\) and \(-1\) eigenspaces of the fermionic \(g\) are separately preserved.  Two real anticommuting matrices that each square to \(-I\) require four real dimensions in either eigenspace.  The two eigenspaces have equal size, so the smallest Nambu representation has \(d=8\).

The tables separate two roles that are useful when designing a Hubbard--Stratonovich decomposition.  Terms in \(A_{\rm group}\) preserve the metric and may fluctuate with either sign.  Terms in \(A_{\rm cone}\) provide damping when their full coefficient matrix obeys Eq.~\eqref{eq:full_coefficient_inequality}.  Since the first Pauli factor fixes whether a term conserves particle number, every class can be read directly as an allowed mixture of hopping, flavor conversion, and anomalous propagation.  The two-dimensional representatives give the most elementary single-flavor structures.  The four- and eight-dimensional representatives supply the internal degrees of freedom required by the corresponding squares and commutation relations.

\section{Group-only classes and time-reversal positivity}
\label{sec:group_only_TR}

When either screening lemma fails, the linear \(J_2\)-odd subspace may still contain matrices, but its intersection with \(\eta A_{\rm cone}\leq0\) contains only zero.  The remaining generators satisfy
\begin{equation}
\begin{gathered}
    gA_{\rm group}+A_{\rm group}^\trans g=0,\qquad
    J_1^\trans A_{\rm group}J_1=\overline{A_{\rm group}},\qquad
    J_2^\trans A_{\rm group}J_2=\overline{A_{\rm group}}.
\end{gathered}
    \label{eq:group_only_generators}
\end{equation}
These are the \(24\) gray cells in each half of Table~\ref{tab:periodic_table}.

Some group-only cells remain sign-problem-free through time-reversal symmetry.  Define anti-linear transformations on the Nambu operators by
\begin{equation}
    \Theta_j=U_j\mc K,\qquad
    U_j=\begin{cases}
        J_j,&s_j=+1,\\
        \ii J_j,&s_j=-1,
    \end{cases}
    \label{eq:induced_time_reversals}
\end{equation}
where \(\mc K\) conjugates numerical coefficients.  The phase in \(U_j\) compensates the sign with which \(J_j\) acts on \(g\).  Consequently,
\begin{equation}
    U_j^\trans gU_j=g,\qquad
    U_j^{-1}\overline{A_{\rm group}}U_j=A_{\rm group}.
    \label{eq:TR_canonical_and_invariant}
\end{equation}
The first equation preserves the canonical commutator or anticommutator, and the second makes every maximal-subgroup slice invariant.  Direct composition gives
\begin{equation}
    \Theta_j^2=q_jI,\qquad
    \Theta_1\Theta_2=(rs_1s_2)\Theta_2\Theta_1.
    \label{eq:induced_TR_algebra}
\end{equation}
The two time reversals anticommute when \(r=-s_1s_2\).  Opposite squares give the Majorana class \(\{+,-\}\), while two negative squares give the Kramers class \(\{-,-\}\)~\cite{li2016prl}.  The time-reversal positivity theorem states that invariance under either pair makes the full bosonic or fermionic Gaussian weight nonnegative, including evolution with anomalous terms.  A derivation for bosons and fermions is given in the companion note \textit{Time-Reversal Symmetries for Sign-Problem-Free Quantum Monte Carlo}.

A physical time reversal must also preserve the conjugacy between creation and annihilation operators.  In the Nambu basis this relation is expressed by the exchange matrix
\begin{equation}
    C=X\otimes I_{N},\qquad C\overline{U_j}C=U_j.
    \label{eq:physical_TR_adjoint}
\end{equation}
For fermions, \(g_{\rm F}=C\), so Eq.~\eqref{eq:physical_TR_adjoint} follows from Eq.~\eqref{eq:TR_canonical_and_invariant} and the unitarity of \(U_j\).  Bosons require this relation in addition to the five signs.  Table~\ref{tab:physical_TR_representatives} supplies a physical representative for every group-only cell marked in Table~\ref{tab:periodic_table}.  Thus eight bosonic cells are time-reversal positive, four in the Majorana class and four in the Kramers class.  Four fermionic cells belong to the Majorana class.

\begin{table}[tb]
\caption{Physical representatives for the group-only classes marked ``TR positive'' in Table~\ref{tab:periodic_table}.  The tuple denotes \((s_1,q_1;s_2,q_2;r)\), and \(d\) is the Nambu dimension.  Pauli words follow Eq.~\eqref{eq:real_pauli_basis}.  For \(s_j=-1\), the physical time reversal uses \(U_j=\ii J_j\).  Every displayed pair satisfies Eq.~\eqref{eq:physical_TR_adjoint}.}
\label{tab:physical_TR_representatives}
\centering
\small
\renewcommand{\arraystretch}{1.15}
\begin{ruledtabular}
\begin{tabular}{ccccc}
Statistics & \((s_1,q_1;s_2,q_2;r)\) & \(d\) & \((J_1,J_2)\) & TR class\\
\hline
\multirow{8}{*}{boson}
& \((+,+;+,-;-)\) & 4 & \((IX,IY)\) & \(\{+,-\}\)\\
& \((+,+;-,-;+)\) & 4 & \((II,ZY)\) & \(\{+,-\}\)\\
& \((+,-;+,-;-)\) & 8 & \((IXY,IYI)\) & \(\{-,-\}\)\\
& \((+,-;-,-;+)\) & 4 & \((IY,ZY)\) & \(\{-,-\}\)\\
& \((-,+;+,-;+)\) & 4 & \((ZI,IY)\) & \(\{+,-\}\)\\
& \((-,+;-,-;-)\) & 4 & \((ZX,ZY)\) & \(\{+,-\}\)\\
& \((-,-;+,-;+)\) & 4 & \((ZY,IY)\) & \(\{-,-\}\)\\
& \((-,-;-,-;-)\) & 8 & \((ZXY,ZYI)\) & \(\{-,-\}\)\\
\hline
\multirow{4}{*}{fermion}
& \((+,-;+,+;-)\) & 4 & \((IY,IX)\) & \(\{+,-\}\)\\
& \((+,-;-,+;+)\) & 4 & \((IY,YY)\) & \(\{+,-\}\)\\
& \((-,-;+,+;+)\) & 2 & \((Y,I)\) & \(\{+,-\}\)\\
& \((-,-;-,+;-)\) & 2 & \((Y,Z)\) & \(\{+,-\}\)\\
\end{tabular}
\end{ruledtabular}
\end{table}

Time-reversal and contraction-semigroup positivity give independent sufficient conditions.  The periodic table marks time reversal only in group-only cells, where it supplies the nonnegative-weight conclusion unavailable from contraction alone.

\section{Summary}

Two real orthogonal matrices organize a complete contraction-semigroup screen for bosonic and fermionic Gaussian quantum Monte Carlo.  The five relations record how each matrix acts on the canonical matrix \(g\), how each squares, and how the pair commutes.  A negative-definite contraction direction exists exactly when \(J_2^2=+I\) for bosons or \(J_2^2=-I\) for fermions, together with \(r=-s_1s_2\).  This gives eight bosonic and eight fermionic semigroup classes.  Four in each set use \(J_2g\), four use \(\ii J_2g\), and the universal negative-definite direction is \(A_{\rm cone}=-\eta\).

The explicit generator tables convert the classification into a construction principle.  Maximal-subgroup directions supply reversible hopping, conversion, and anomalous Hubbard--Stratonovich channels.  Cone directions supply the damping that prevents determinant zeros and bosonic poles.  When the auxiliary-field domain is connected to a positive reference, and a bosonic realization has a convergent many-body trace, these directions define a sign-free semigroup construction.  The remaining group-only cells admit a second screen based on time-reversal symmetry, adding eight bosonic and four fermionic sign-free classes.  Together, the two routes give a systematic map from quadratic symmetry data to sign-problem-free interacting models.

\bibliography{references}

\end{document}
